\documentclass[journal]{IEEEtran}

\usepackage{cite}
\usepackage{amsmath,amssymb,amsfonts}
\usepackage{graphicx}
\usepackage{textcomp}
\usepackage{xcolor}
\usepackage{booktabs}
\usepackage{multirow}
\usepackage[utf8]{inputenc}
\usepackage[T1]{fontenc}
\usepackage{microtype}
\usepackage{array}
\usepackage{algorithm}
\usepackage{algorithmic}
\usepackage{amsthm}
\newtheorem{assumption}{Assumption}
\newtheorem{theorem}{Theorem}
\newtheorem{proposition}{Proposition}
\newtheorem{remark}{Remark}

\interdisplaylinepenalty=2500
\emergencystretch=2em
\hyphenation{Koop-man trajec-tory trajec-tories auto-regressive regu-lar-ization data-set data-sets non-over-lapping pre-train-ing Trans-former state-em-bedded delay-em-bed-ding}

\def\BibTeX{{\rm B\kern-.05em{\sc i\kern-.025em b}\kern-.08em
    T\kern-.1667em\lower.7ex\hbox{E}\kern-.125emX}}

\begin{document}

\title{Transformer-Koopman: Finite-History Lifting for Controlled Nonlinear Systems}

\author{Bolin Tao}

\markboth{IEEE Journal Draft}{Tao: Transformer-Koopman}

\maketitle

\begin{abstract}
Koopman modeling is attractive for controlled nonlinear systems because it seeks a lifted representation in which nonlinear dynamics become approximately linear. In practice, however, the quality of the resulting model depends critically on the lifting map. Existing learned Koopman models often construct observables from the instantaneous state alone, which can discard short-term temporal information useful for identifying accurate linear latent dynamics. This paper proposes Transformer-Koopman, a temporal lifting framework that learns Koopman observables from finite history windows by combining non-overlapping multivariate temporal patches with a Transformer encoder. The lifted state evolves according to linear latent dynamics controlled by the input, and the associated system matrices are identified through a hybrid pipeline that combines encoder pretraining with Extended Dynamic Mode Decomposition (EDMD) style regression. Open-loop experiments on a 6-state manipulator dataset show substantially lower rollout error than traditional EDMD and multi-layer perceptron (MLP) Koopman baselines, while ablation studies on a second 2-state dataset demonstrate that history length, patch design, and moderate regularization materially affect model quality. The resulting representation remains compact, interpretable, and naturally compatible with subsequent tracking-control studies.
\end{abstract}

\begin{IEEEkeywords}
Koopman operator, controlled nonlinear systems, temporal lifting, Transformer, data-driven modeling, trajectory prediction
\end{IEEEkeywords}

\section{Introduction}
\label{sec:introduction}

Learning predictive models for nonlinear controlled systems is a central challenge in system identification and model-based control \cite{ljung1999system, vanoverschee1994n4sid, nelles2020nonlinear}. For control-oriented applications, model quality is not measured only by one-step fit, but also by whether the representation supports reliable multi-step prediction and a clean interface for downstream controller design. The Koopman framework is attractive in this regard: it seeks a lifted coordinate space in which nonlinear state evolution is governed, at least approximately, by linear dynamics \cite{koopman1931hamiltonian, brunton2022modern}. When such a representation exists, standard linear tools for prediction and control become directly applicable even though the underlying physical system remains nonlinear \cite{otto2021koopman, korda2018linear}.

Despite this appeal, constructing an accurate and compact lifting map remains the central open problem. Classical data-driven Koopman methods rely on hand-crafted observable dictionaries \cite{williams2015data, klus2018data}, and learned lifting approaches \cite{lusch2018deep, takeishi2017koopman} have improved flexibility considerably, but most still encode the lifted state from the instantaneous physical state alone. This is a modeling limitation: for many nonlinear systems, a short window of recent trajectory context carries information relevant to constructing a more accurate linear latent evolution law, yet this context is discarded when the lifting depends only on the current state. The present paper addresses this gap by learning the lifting map from finite history windows rather than from the instantaneous state alone.

Rather than mapping the current state to the lifted space, the proposed method encodes a short finite-history window using a compact Transformer encoder \cite{vaswani2017attention} with non-overlapping temporal patch tokenization \cite{nie2023time}. The resulting history-dependent feature vector is appended to the current physical state to form the lifted state, which is then propagated by an identified linear controlled dynamical model. Because the physical state is embedded explicitly in the lifted vector, exact state recovery is achieved through a fixed linear projection, avoiding the need for a learned nonlinear decoder and preserving a clean interface for subsequent controller design.

The proposed framework, called \emph{Transformer-Koopman}, is trained through a three-stage pipeline. The encoder and linear dynamics are first jointly optimized end-to-end via gradient descent, then the linear system matrices are refit in closed form using Extended Dynamic Mode Decomposition (EDMD) under the converged encoder \cite{williams2015data, proctor2018generalizing}. This hybrid procedure combines the representational flexibility of learned temporal lifting with the favorable conditioning properties of closed-form linear regression. Recent evidence that compact Transformer encoders can capture meaningful dynamical structure \cite{forootani2026compact, forootani2025synthetic} further supports the viability of this design.

The contributions of this paper are as follows.
\begin{enumerate}
    \item We propose a temporal lifting framework that constructs Koopman observables from finite history windows using non-overlapping temporal patches and a Transformer encoder, rather than from the instantaneous state alone.
    \item We construct a compact lifted model in which the physical state is recovered by a fixed projection, preserving interpretability and a clean control interface without requiring a learned nonlinear decoder.
    \item We present a three-stage hybrid training pipeline that jointly optimizes the encoder and linear dynamics via gradient descent, followed by a closed-form EDMD refit, and validate the approach through open-loop trajectory-prediction comparison and ablation analysis.
\end{enumerate}

The rest of the paper is organized as follows. Section~\ref{sec:related_work} reviews related work. Section~\ref{sec:problem} formulates the learning problem. Section~\ref{sec:methodology} presents the proposed method. Section~\ref{sec:control} presents the controller design and stability analysis. Section~\ref{sec:experiments} reports experiments. Section~\ref{sec:conclusion} concludes.

\section{Related Work}
\label{sec:related_work}

\subsection{Koopman Modeling: Foundations, Learning, and Control}

Koopman operator theory recasts nonlinear dynamics as a linear evolution law acting on functions of the state, called observables \cite{koopman1931hamiltonian, koopman1932spectra, mezic2005spectral, rowley2009spectral}. In data-driven settings, Extended Dynamic Mode Decomposition (EDMD) and related transfer-operator methods approximate this operator from trajectory data by fitting a linear predictor over a finite dictionary of observables \cite{williams2015data, williams2015kernel, klus2018data}. These methods offer interpretable spectral structure and closed-form regression, but their accuracy depends critically on the expressiveness of the hand-crafted dictionary.

Learned lifting approaches remove this dependence by discovering the observable map from data. Representative architectures include invariant-subspace encoders \cite{takeishi2017koopman}, deep Koopman autoencoders \cite{lusch2018deep, yeung2019learning}, variational formulations \cite{morton2019deep, yuan2022variational}, and consistency-regularized autoencoders \cite{azencot2020forecasting}. A related line of work emphasizes temporal structure: linearly recurrent autoencoders \cite{otto2019linearly} and recurrent Koopman variants \cite{rank2019recurrent} show that incorporating sequential context can improve the quality of the learned linearization. Data-driven models such as DynoNet \cite{forgione2021dynonet} and SINDy \cite{brunton2016sindy} further motivate combining structural priors with flexible function approximation for dynamical systems.

When control inputs are present, the Koopman framework is particularly valuable because linear latent dynamics are directly compatible with standard control synthesis. DMDc \cite{proctor2016dmdc}, input-aware Koopman formulations \cite{proctor2018generalizing, brunton2016invariant}, and Koopman-based MPC \cite{korda2018linear} established this direction. More recent work extends it to learned representations \cite{han2020deep, shi2022deep} and provides closed-loop guarantees and error bounds \cite{strasser2026overview, duran2025control}. Across this body of work, a consistent finding is that controller performance is tightly coupled to model quality, which motivates the emphasis of the present paper on improving the lifting map itself.

\subsection{Transformers for Temporal Sequence Modeling}

The Transformer architecture \cite{vaswani2017attention} has become a standard building block for sequence modeling, with self-attention providing a flexible mechanism for capturing long-range dependencies. Time-series adaptations include efficient-attention variants \cite{zhou2021informer, wang2020linformer, katharopoulos2020transformers, choromanski2021performer}, decomposition-based models \cite{wu2021autoformer, zhou2022fedformer}, and reformulations that invert the channel-time structure \cite{zeng2023transformers, liu2024itransformer}. Patch-wise tokenization \cite{nie2023time} has proven particularly effective: partitioning the input sequence into non-overlapping segments before attention reduces token count while preserving local temporal structure. Compact Transformer studies demonstrate that lightweight encoders already capture meaningful dynamical structure under controlled synthetic conditions \cite{forootani2026compact, forootani2025synthetic}.

Recent hybrid architectures have begun combining Transformer representations with Koopman-style latent dynamics for time-series forecasting \cite{forootani2026deepkoopformer, forootani2026spectral}. These works are the most closely related to ours in architecture, but differ in objective: they target generic sequence forecasting, whereas the present work targets controlled nonlinear system identification. Concretely, our method preserves the physical state explicitly in the lifted vector, identifies the linear dynamics via a hybrid gradient-EDMD pipeline rather than end-to-end forecasting loss alone, and is designed to provide a clean interface for subsequent controller design rather than to minimize multi-step forecast error on held-out time series.

\subsection{Positioning of the Present Work}

The two lines of research reviewed above — learned Koopman lifting and compact Transformer encoders — have largely developed in parallel. The present paper connects them in a control-oriented setting. Unlike forecasting-oriented hybrid models, Transformer-Koopman is designed so that the Transformer serves only as a history compression module: it constructs the lifting map, while all prediction is performed by the identified linear latent dynamics. This separation keeps the model compact, makes the control interface explicit, and allows the linear system matrices to be identified by closed-form regression rather than absorbed into a black-box sequence model. In this sense, Transformer-Koopman bridges the gap between learned temporal sequence modeling and control-oriented Koopman identification in a way that neither line of prior work has directly addressed.

\section{Preliminaries and Problem Formulation}

\subsection{Preliminary: Koopman Modeling and Temporal Lifting}
\label{sec:preliminary}

The Koopman operator provides a linear representation of nonlinear dynamics by acting on functions of the state, called observables, rather than on the state itself \cite{koopman1931hamiltonian, mezic2005spectral, rowley2009spectral, brunton2022modern}. For a discrete-time controlled system $x_{k+1} = F(x_k,u_k)$, the Koopman evolution of an observable $\varphi:\mathbb{R}^n\to\mathbb{R}$ is
\begin{equation}
    (\mathcal{K}\varphi)(x) = \varphi(F(x,u)).
    \label{eq:koopman_operator}
\end{equation}
Intuitively, if suitable observables are chosen, the nonlinear evolution of $x_k$ may appear linear in the lifted space, allowing one to leverage linear techniques for analysis and control.

Since the Koopman operator is infinite-dimensional, practical computation requires a finite-dimensional approximation. Extended Dynamic Mode Decomposition (EDMD) achieves this by selecting a dictionary of $d$ observables $g_1(x),\ldots,g_d(x)$, stacking them into
\begin{equation}
    z_k = g(x_k) \in \mathbb{R}^d,
    \label{eq:instantaneous_lifting}
\end{equation}
and solving for matrices $A$ and $B$ that minimize the prediction error
\begin{equation}
    z_{k+1} \approx A z_k + B u_k.
    \label{eq:learned_dynamics}
\end{equation}
EDMD preserves linear structure and allows spectral analysis, but its performance depends critically on the chosen dictionary and may fail to capture short-term dependencies.

To address this, temporal lifting incorporates a short history of past states. For a history length $P$, define
\begin{equation}
    X_k = [x_{k-P+1}^\top, \ldots, x_k^\top]^\top \in \mathbb{R}^{P\times n}.
    \label{eq:history_window}
\end{equation}
A temporal lifting map $g_\theta$ then produces
\begin{equation}
    z_k = g_\theta(X_k),
    \label{eq:temporal_lifting}
\end{equation}
which encodes both the current state and recent trajectory context. This approach leverages ideas from delay-coordinate embeddings \cite{arbabi2017ergodic, rank2019recurrent}, allowing the linear latent dynamics
\begin{equation}
    z_{k+1} = A z_k + B u_k + w_k
\end{equation}
to more accurately approximate the true nonlinear system, where $w_k$ is the residual error.

In this work, $g_\theta$ is implemented with a Transformer encoder that compresses the history window into a compact feature vector appended to the current state. This architecture is described in detail in Section~\ref{sec:temporal_lifting}.


\subsection{Problem Formulation}
\label{sec:problem}

Consider a discrete-time controlled nonlinear system
\begin{equation}
    x_{k+1} = F(x_k,u_k),
    \label{eq:nonlinear_system}
\end{equation}
where $x_k \in \mathbb{R}^n$ is the state, $u_k \in \mathbb{R}^m$ is the control input, and $F(\cdot)$ is an unknown nonlinear transition map. We assume that full-state measurements and the corresponding control inputs are available from trajectory data. This assumption matches the datasets and learning pipeline considered in this study.

The history window $X_k$ is defined in \eqref{eq:history_window}. From recorded trajectories, we form samples of the form
\begin{equation}
    \mathcal{D} = \{(X_k,u_k,x_{k+1})\}_{k=1}^{M},
    \label{eq:dataset}
\end{equation}
where only windows that remain within a single trajectory are retained. In the actual dataset construction, both the history window and the one-step prediction target are required to respect trajectory boundaries.

The modeling goal is to learn a finite-dimensional lifted representation in which the controlled nonlinear dynamics can be approximated by a linear latent system. To this end, we define a temporal lifting map $z_k = g_\theta(X_k)$ \emph{as in} \eqref{eq:temporal_lifting}, where $z_k \in \mathbb{R}^d$ is a lifted state derived from a finite history window rather than from the instantaneous state alone. The target lifted dynamics follow $z_{k+1}=Az_k+Bu_k$ (cf. \eqref{eq:instantaneous_lifting}) with residual $w_k$:
\begin{equation}
    z_{k+1} = A z_k + B u_k + w_k,
    \label{eq:lifted_dynamics}
\end{equation}
where $A \in \mathbb{R}^{d\times d}$ and $B \in \mathbb{R}^{d\times m}$ are identified matrices and $w_k$ denotes the finite-dimensional approximation residual.

This formulation reflects the central hypothesis of the paper. If recent history carries information that is useful for constructing an approximately linear latent evolution law, then a finite-history lifting map can produce a more accurate Koopman model than an instantaneous-state encoder. At the same time, we do not simply treat the stacked history itself as the lifted state. Instead, we seek a compact lifted representation that preserves the current physical state explicitly while using the remaining latent coordinates to encode history-dependent information.

Two practical remarks are important. First, the present encoder uses \emph{state history only}; past controls are not included in the encoder input. This design keeps the lifting map simple and matches the training and data-construction procedure used here, and it should be understood as a modeling choice specific to the present formulation rather than as a general requirement of temporal lifting. Second, the primary objective of this paper is predictive model quality under autoregressive rollout, with control compatibility treated as a downstream structural requirement rather than as a completed closed-loop claim.

% =========================================================
\section{Transformer-Based Temporal Lifting}
\label{sec:methodology}
% =========================================================

This section presents the proposed model and its training procedure from the perspective of control-compatible model identification.

% ---------------------------------------------------------
\subsection{Temporal Lifting Architecture}
\label{sec:temporal_lifting}
% ---------------------------------------------------------

The encoder consists of three stages: temporal patch tokenization, Transformer encoding, and lifted state construction.

The encoder takes the history window $X_k \in \mathbb{R}^{P\times n}$ as input. Assuming $p$ divides $P$, the window is partitioned into $N=P/p$ contiguous non-overlapping patches of temporal length $p$. Denote the $i$th patch by
\begin{equation}
    X_k^{(i)} \in \mathbb{R}^{p\times n}, \qquad i=1,\ldots,N.
\end{equation}
Each patch contains all $n$ state channels over $p$ consecutive time steps, so the tokenization is \emph{jointly multivariate} rather than channel-wise. After vectorization,
\begin{equation}
    q_k^{(i)} = \mathrm{vec}(X_k^{(i)}) \in \mathbb{R}^{pn},
\end{equation}
the corresponding token embedding is
\begin{equation}
    e_k^{(i)} = W_p q_k^{(i)} + b_p,
    \label{eq:patch_embedding}
\end{equation}
where $W_p \in \mathbb{R}^{d_{\mathrm{model}}\times pn}$ and $b_p \in \mathbb{R}^{d_{\mathrm{model}}}$ are learnable parameters. Here $d_{\mathrm{model}}$ denotes the Transformer embedding dimension.

This tokenization performs two functions simultaneously. It aggregates short local temporal patterns before global sequence modeling, and it reduces the effective token length from $P$ to $P/p$. As a result, the self-attention stage operates on a shorter sequence, which is useful when the objective is to learn a compact temporal lifting module rather than to maximize raw sequence length.

The patch embeddings are augmented with sinusoidal positional encoding and processed by a Transformer encoder:
\begin{equation}
    \tilde{e}_k^{(i)} = e_k^{(i)} + s^{(i)},
\end{equation}
where $s^{(i)}$ denotes the fixed positional encoding associated with token index $i$. Let
\begin{equation}
    H_k = \mathrm{Transformer}(\tilde{e}_k^{(1)},\ldots,\tilde{e}_k^{(N)})
\end{equation}
be the encoded token sequence. The implementation then applies a pooling-based readout operator,
\begin{equation}
    r_k = \mathrm{Pool}(H_k),
\end{equation}
where the pooling rule can be mean, max, or last-token selection, followed by a linear projection
\begin{equation}
    h_k = W_h r_k + b_h \in \mathbb{R}^{d-n}.
    \label{eq:history_feature}
\end{equation}
The lifted state is constructed as
\begin{equation}
    z_k = [x_k;h_k] \in \mathbb{R}^{d},
    \label{eq:lifted_state_form}
\end{equation}
where $x_k$ is the final state in the history window and $h_k$ is the learned temporal feature. The first $n$ coordinates coincide with the physical state. The fixed projection $x_k=Cz_k$ with $C=[I_n\;0]$ recovers $x_k$ exactly:
\begin{equation}
    x_k = C z_k, \qquad C = [I_n\;0].
    \label{eq:decoder_projection}
\end{equation}

Given this lifted state, the latent dynamics follow
\begin{equation}
    z_{k+1} = A z_k + B u_k,
    \label{eq:temporal_dynamics}
\end{equation}
where $A$ and $B$ are identified in the training procedure described next. The residual $w_k$ from the problem formulation is absorbed into the learned encoder residual in this deterministic model statement.

% ---------------------------------------------------------
\subsection{Training Procedure}
\label{sec:training}
% ---------------------------------------------------------

The training pipeline consists of three stages that combine neural temporal lifting with closed-form linear regression in the learned latent space, and is summarized in Algorithm~\ref{alg:training}.

\textbf{Stage 0: Initial EDMD estimate.}
With the encoder at its random initialization, all training windows are mapped to latent states and an initial estimate of $(A,B)$ is obtained via the regularized least-squares regression defined in Stage~2. This warm-start provides a meaningful linearization point before joint gradient-based optimization begins.

\textbf{Stage 1: Joint pretraining.}
Let the shifted next-window be
\begin{equation}
    X_k^{+} = [x_{k-P+2}^\top,\,\ldots,\,x_{k+1}^\top]^\top,
\end{equation}
and define the true next latent state as $z_k^{+} = g_\theta(X_k^{+})$.
The predicted next latent and physical states are
\begin{equation}
    \hat{z}_{k+1} = Az_k + Bu_k,
    \qquad
    \hat{x}_{k+1} = C\hat{z}_{k+1}.
\end{equation}
Three losses drive the joint pretraining. The physical-state prediction loss penalizes rollout error directly in the observable space:
\begin{equation}
    \mathcal{L}_{x} = \|\hat{x}_{k+1}-x_{k+1}\|_2^2.
    \label{eq:loss_x}
\end{equation}
The latent-state prediction loss encourages the linear propagation $Az_k+Bu_k$ to match the encoder output at the next history window, aligning the learned dynamics with the lifting map:
\begin{equation}
    \mathcal{L}_{z} = \|\hat{z}_{k+1}-z_k^{+}\|_2^2.
    \label{eq:loss_z}
\end{equation}
The latent consistency loss imposes a self-consistency constraint on the predicted trajectory: the history window constructed by replacing the most recent entry with the predicted next state $\hat{x}_{k+1}$ is re-encoded, and the result is required to agree with the true next latent state,
\begin{equation}
    \mathcal{L}_{c} = \bigl\|g_\theta\!\left(\mathcal{W}^{+}_{\mathrm{pred}}\right)-z_k^{+}\bigr\|_2^2,
    \label{eq:loss_c}
\end{equation}
where $\mathcal{W}^{+}_{\mathrm{pred}}$ denotes the shifted history window whose most recent entry is $\hat{x}_{k+1}$. This term penalizes encoder inconsistencies that would accumulate during autoregressive rollout, since at each rollout step the encoder is re-evaluated on a window containing previously predicted states. The overall pretraining objective is
\begin{equation}
    \mathcal{L}_{\mathrm{pre}} = \lambda_x\mathcal{L}_x + \lambda_z\mathcal{L}_z + \lambda_c\mathcal{L}_c.
    \label{eq:pretraining_loss}
\end{equation}
Both the encoder parameters $\theta$ and the Koopman matrices $(A,B)$ are treated as learnable parameters and updated jointly by gradient descent on $\mathcal{L}_{\mathrm{pre}}$. This joint optimization allows the lifting map and the linear dynamics to co-adapt, so that the encoder learns observables that are amenable to linear propagation while $(A,B)$ simultaneously adjusts to the evolving latent geometry.

\textbf{Stage 2: Final EDMD refit.}
After joint pretraining, the encoder parameters $\theta$ are frozen and latent trajectories are recomputed with the converged $\theta$. The linear dynamics are then refit in closed form via EDMD, replacing the gradient-trained $(A,B)$ with the globally optimal least-squares solution under the fixed encoder. Collect the latent states, shifted latent states, and control inputs into matrices
\begin{align}
    Z &= [z_1,\ldots,z_M]^\top, \\
    Z^{+} &= [z_1^{+},\ldots,z_M^{+}]^\top, \\
    U &= [u_1,\ldots,u_M]^\top,
\end{align}
and define the augmented regressor $\Omega=[Z\;U]$. The ridge-regularized EDMD estimate is then
\begin{equation}
    \Theta = (Z^{+})^\top \Omega
    \left(\Omega^\top\Omega + \lambda I\right)^{-1},
    \label{eq:edmd_solution}
\end{equation}
where $\Theta \in \mathbb{R}^{d\times(d+m)}$ is partitioned as $\Theta = [A\;B]$. The regularization parameter $\lambda > 0$ controls ill-conditioning in the latent regression; its effect on model quality is examined in the ablation study of Section~\ref{sec:experiments}.

The complete pipeline is summarized in Algorithm~\ref{alg:training}.

\begin{algorithm}[!t]
\caption{Transformer-Koopman Training Pipeline}
\label{alg:training}
\begin{algorithmic}[1]
\REQUIRE Dataset $\mathcal{D}$ \eqref{eq:dataset}; hyperparameters $P,p,d,\lambda_x,\lambda_z,\lambda_c,\lambda,E$.
\ENSURE Encoder parameters $\theta^{*}$; system matrices $A^{*},B^{*}$.

\STATE \textbf{// Stage 0: Initial EDMD estimate}
\STATE Initialize $\theta$ randomly; encode all windows to obtain $Z,Z^{+},U$.
\STATE $(A,B)\leftarrow$ solve \eqref{eq:edmd_solution}. \hfill\COMMENT{warm-start for Stage 1}

\STATE \textbf{// Stage 1: Joint pretraining of $\theta$, $A$, $B$}
\FOR{epoch $e=1$ \TO $E$}
    \FOR{each mini-batch from $\mathcal{D}$}
        \STATE Compute $\hat{z}_{k+1}=Az_k+Bu_k$,\; $\hat{x}_{k+1}=C\hat{z}_{k+1}$.
        \STATE $(\theta,A,B)\leftarrow(\theta,A,B) - \eta\,\nabla_{\theta,A,B}\mathcal{L}_{\mathrm{pre}}$ \eqref{eq:pretraining_loss}.
    \ENDFOR
\ENDFOR

\STATE \textbf{// Stage 2: Final EDMD refit (encoder frozen)}
\STATE Fix $\theta$; re-encode all windows to obtain updated $Z,Z^{+},U$.
\STATE $(A^{*},B^{*})\leftarrow$ solve \eqref{eq:edmd_solution}.
\RETURN $\theta^{*}\leftarrow\theta$,\;$A^{*}$,\;$B^{*}$.
\end{algorithmic}
\end{algorithm}


\subsection{Modeling Error Characterization}
\label{subsec:modeling_error}

The lifted dynamics identified above approximate the true system evolution. In practice, the finite-dimensional Koopman approximation and the learned encoder both introduce a residual. Define the one-step modeling error as
\begin{equation}
    w_k = z_{k+1} - (Az_k + Bu_k),
    \label{eq:modeling_error}
\end{equation}
where $z_{k+1} = g_\theta(X_{k+1})$ is the true next lifted state obtained by re-encoding the next history window. The error $w_k$ accumulates contributions from two sources: the finite-dimensional truncation of the Koopman operator, and the encoder approximation residual that persists after training. Under the joint training objective of Section~\ref{sec:training}, both contributions are minimized, but not driven to zero in general.

Similarly, the lifted reference $z_k^{\mathrm{ref}} = g_\theta(X_k^{\mathrm{ref}})$ satisfies
\begin{equation}
    z_{k+1}^{\mathrm{ref}} = Az_k^{\mathrm{ref}} + Bu_k^{\mathrm{ref}} + w_k^{\mathrm{ref}},
    \label{eq:ref_error}
\end{equation}
where $w_k^{\mathrm{ref}}$ is the corresponding reference modeling error. Both $w_k$ and $w_k^{\mathrm{ref}}$ arise from the same encoder and operator approximation, so they share the same character.


% ---------------------------------------------------------
\subsection{Autoregressive Multi-Step Rollout}
\label{sec:rollout}
% ---------------------------------------------------------

After training, the model is deployed autoregressively for open-loop trajectory prediction. Starting from an initial history window $X_0 \in \mathbb{R}^{P \times n}$ and a future control sequence $\{u_0,\ldots,u_{T-1}\}$, the model iterates the following steps for $t=0,1,\ldots,T-1$:
\begin{equation}
    z_t = g_{\theta^*}(\mathcal{W}_t), \quad
    \hat{z}_{t+1} = A^* z_t + B^* u_t, \quad
    \hat{x}_{t+1} = C\,\hat{z}_{t+1},
    \label{eq:rollout_step}
\end{equation}
where $\mathcal{W}_t$ denotes the current sliding history window, which is updated at each step by appending $\hat{x}_{t+1}$ and discarding the oldest entry.

Rather than propagating a single latent state indefinitely through the linear operator, the model re-evaluates the history-based lifting map at every step. This design confines error accumulation in the latent propagation to a single linear step per iteration, while the encoder continuously re-anchors the latent state to a recent portion of the predicted trajectory. The rollout produces the predicted state sequence $\{\hat{x}_1,\hat{x}_2,\ldots,\hat{x}_T\}$, which is used directly for evaluation in Section~\ref{sec:experiments}.

% =========================================================
\section{Controller Design and Stability Analysis}
\label{sec:control}
% =========================================================

\subsection{Controller Design}
\label{subsec:controller_design}

The linear lifted dynamics identified in Section~\ref{sec:methodology} are directly compatible with standard linear control synthesis. The lifted state has the explicit block structure
\begin{equation}
    z_k = [x_k;\, h_k] \in \mathbb{R}^d,
    \label{eq:lifted_block}
\end{equation}
where $x_k \in \mathbb{R}^n$ occupies the first $n$ coordinates (the physical state, copied directly into the lifted vector) and $h_k \in \mathbb{R}^{d-n}$ contains the history-dependent features learned by the Transformer encoder. This structure motivates a hierarchical control design that assigns different priorities to the two components of the lifted error.

\subsubsection*{Reference trajectory and structured error decomposition}

Assume that a reference physical trajectory $x_k^{\mathrm{ref}}$ and reference input sequence $u_k^{\mathrm{ref}}$ are available. The corresponding reference lifted state is formed as
\begin{equation}
    z_k^{\mathrm{ref}} = [x_k^{\mathrm{ref}};\, h_k^{\mathrm{ref}}],
    \label{eq:ref_lifted}
\end{equation}
where $h_k^{\mathrm{ref}} = W_h\,\mathrm{Pool}\!\left(\mathrm{Transformer}(\tilde{e}^{(1:N),\mathrm{ref}}_k)\right) + b_h$ is the encoder output on the reference history window $X_k^{\mathrm{ref}} = [(x_{k-P+1}^{\mathrm{ref}})^\top,\ldots,(x_k^{\mathrm{ref}})^\top]^\top$.

The lifted tracking error $e_k^z = z_k - z_k^{\mathrm{ref}}$ inherits the block structure of \eqref{eq:lifted_block}:
\begin{equation}
    e_k^z = \begin{bmatrix} e_k^x \\ e_k^h \end{bmatrix}
           = \begin{bmatrix} x_k - x_k^{\mathrm{ref}} \\ h_k - h_k^{\mathrm{ref}} \end{bmatrix},
    \label{eq:structured_error}
\end{equation}
where $e_k^x \in \mathbb{R}^n$ is the physical-state error and $e_k^h \in \mathbb{R}^{d-n}$ is the memory-state error. Because the physical state is recovered exactly via $x_k = Cz_k$ with $C=[I_n\;0]$, one has
\begin{equation}
    e_k^x = x_k - x_k^{\mathrm{ref}} = C\,e_k^z.
    \label{eq:phys_error_proj}
\end{equation}

\subsubsection*{Hierarchical control law}

The total control input is decomposed into a feedforward term and a hierarchical feedback term:
\begin{equation}
    u_k = u_{\mathrm{ff},k} + u_{\mathrm{fb},k}.
    \label{eq:ctrl_decomp}
\end{equation}
The feedforward term compensates for the reference trajectory in the lifted space:
\begin{equation}
    u_{\mathrm{ff},k} = B^\dagger\!\left(z_{k+1}^{\mathrm{ref}} - A z_k^{\mathrm{ref}}\right),
    \label{eq:feedforward}
\end{equation}
where $B^\dagger = (B^\top B)^{-1}B^\top$ is the Moore--Penrose pseudoinverse of $B$. The hierarchical feedback term penalizes the structured error \eqref{eq:structured_error} according to the two-level priority:
\begin{equation}
    u_{\mathrm{fb},k} = -K_x e_k^x - K_h e_k^h = -K_c\, e_k^z,
    \label{eq:hierarchical_fb}
\end{equation}
where $K_c = [K_x\;|\;K_h]$ is partitioned conformably with \eqref{eq:structured_error}. The complete hierarchical control law is therefore
\begin{equation}
    u_k = B^\dagger\!\left(z_{k+1}^{\mathrm{ref}} - A z_k^{\mathrm{ref}}\right) - K_x e_k^x - K_h e_k^h.
    \label{eq:full_ctrl}
\end{equation}

\subsubsection*{Block-diagonal weight matrix design}

The gain $K_c$ is computed from a single discrete-time algebraic Riccati equation (DARE) for the pair $(A,B)$, with cost matrices $Q$ and $R \succ 0$. The key design choice is the block-diagonal structure of $Q$:
\begin{equation}
    Q = \begin{bmatrix} Q_x & 0 \\ 0 & \varepsilon Q_h \end{bmatrix},
    \qquad \varepsilon \ll 1,
    \label{eq:block_diag_Q}
\end{equation}
where $Q_x \succ 0$ is the physical-state weight matrix, $Q_h \succ 0$ is the memory-state weight matrix, and $\varepsilon \in (0,1)$ is the hierarchical strength factor. The LQR objective is
\begin{equation}
    J = \sum_{k=0}^{\infty}
    \left[(e_k^z)^\top Q\, e_k^z + (\delta u_k)^\top R\,\delta u_k\right],
    \label{eq:lqr_objective}
\end{equation}
with $\delta u_k = u_{\mathrm{fb},k}$. Solving the DARE
\begin{equation}
    A^\top P A - P - A^\top P B(R + B^\top P B)^{-1}B^\top P A + Q = 0
    \label{eq:dare}
\end{equation}
for $P \succ 0$ yields the unified hierarchical gain
\begin{equation}
    K_c = (R + B^\top P B)^{-1}B^\top P A = [K_x\;|\;K_h].
    \label{eq:lqr_gain}
\end{equation}
Because $Q_x \gg \varepsilon Q_h$, the LQR optimization naturally produces $\|K_x\| \gg \|K_h\|$, so the physical-state error dominates the control response while the memory-state error acts as an auxiliary correction without driving the controller. The notation $K_c$ is used to avoid conflict with the regression matrix in \eqref{eq:edmd_solution}.

\begin{remark}
The gains $K_x$ and $K_h$ are obtained from a single Riccati equation rather than designed independently. The hierarchical effect emerges entirely from the weight ratio $Q_x / (\varepsilon Q_h)$: decreasing $\varepsilon$ amplifies physical-state priority without altering the stability guarantee provided by the LQR framework.
\end{remark}

Because physical-state recovery is exact through \eqref{eq:decoder_projection}, the physical tracking error satisfies
\begin{equation}
    x_k - x_k^{\mathrm{ref}} = C\,e_k^z,
\end{equation}
so tracking performance in the original state space can be evaluated directly from the lifted error without a separate observation map. The controller structure above is intentionally kept fixed across all model variants; in the planned trajectory-tracking study described in Section~\ref{sec:experiments}, only the learned Koopman model will be varied, so that differences in closed-loop performance can be attributed to modeling quality rather than controller design.

% ---------------------------------------------------------
\subsection{Stability Analysis}
\label{subsec:stability}
% ---------------------------------------------------------

This subsection presents the stability analysis of the closed-loop tracking system formed by the learned Koopman model and the hierarchical LQR controller derived in Section~\ref{subsec:controller_design}. The analysis proceeds from the modeling error characterization of Section~\ref{subsec:modeling_error}, through a boundedness assumption, to a formal ultimate-boundedness result for the lifted error $e_k^z$.

\subsubsection*{Closed-loop error dynamics}

Substituting the control law \eqref{eq:full_ctrl} into the lifted dynamics \eqref{eq:lifted_dynamics} and subtracting the reference dynamics \eqref{eq:ref_error}, the closed-loop error evolves as
\begin{equation}
    e_{k+1}^z = A_{\mathrm{cl}}\, e_k^z + \delta_k,
    \label{eq:cl_error_dyn}
\end{equation}
where $A_{\mathrm{cl}} = A - BK_c$ is the closed-loop system matrix and $\delta_k = w_k - w_k^{\mathrm{ref}}$ is the net disturbance arising from the modeling residuals in \eqref{eq:modeling_error} and \eqref{eq:ref_error}.

The stability analysis rests on the following assumption.

\begin{assumption}[Uniform error bound]
\label{assump:bound}
There exists a constant $\varepsilon \geq 0$ such that
\begin{equation}
    \|w_k\| \leq \varepsilon \quad \text{and} \quad \|w_k^{\mathrm{ref}}\| \leq \varepsilon
    \label{eq:error_bound}
\end{equation}
for all $k \geq 0$ along the trajectories of interest.
\end{assumption}

\begin{remark}
Assumption~\ref{assump:bound} is standard in approximate Koopman control \cite{strasser2026overview} and is supported empirically by the bounded rollout errors observed in Section~\ref{sec:experiments}. The bound $\varepsilon$ can in principle be estimated from validation data as the maximum one-step residual after training. Under this assumption, $\|\delta_k\| \leq 2\varepsilon$ for all $k \geq 0$.
\end{remark}

\subsubsection*{Ultimate boundedness}

\begin{proposition}[Ultimate boundedness]
\label{prop:ub}
Let $A_{\mathrm{cl}} = A - BK_c$, where $K_c = [K_x\;|\;K_h]$ is the hierarchical LQR gain computed from \eqref{eq:lqr_gain} with $Q$ as in \eqref{eq:block_diag_Q} and $R \succ 0$, and suppose $(A,B)$ is stabilizable and $(A,Q^{1/2})$ is detectable. Then $A_{\mathrm{cl}}$ is Schur stable, i.e., $\rho(A_{\mathrm{cl}}) < 1$. Suppose further that $\|A_{\mathrm{cl}}\| < 1$ in the induced $2$-norm. Under Assumption~\ref{assump:bound}, the lifted tracking error $e_k^z$ is ultimately bounded: there exists a finite $k_0$ such that for all $k \geq k_0$,
\begin{equation}
    \|e_k^z\| \leq \frac{2\varepsilon}{1 - \|A_{\mathrm{cl}}\|}.
    \label{eq:ub_bound}
\end{equation}
\end{proposition}

\begin{proof}[Proof sketch]
Since $(A,B)$ is stabilizable and $(A,Q^{1/2})$ is detectable, the DARE \eqref{eq:dare} has a unique positive-definite solution $P \succ 0$, and the resulting $A_{\mathrm{cl}} = A - BK_c$ satisfies $\rho(A_{\mathrm{cl}}) < 1$ \cite{korda2018linear}. Choose $V_k = (e_k^z)^\top P\, e_k^z$ as a Lyapunov function. From \eqref{eq:cl_error_dyn},
\begin{equation}
    V_{k+1} = (A_{\mathrm{cl}} e_k^z + \delta_k)^\top P (A_{\mathrm{cl}} e_k^z + \delta_k).
\end{equation}
Expanding and applying the LQR optimality condition $A_{\mathrm{cl}}^\top P A_{\mathrm{cl}} - P = -(Q + K_c^\top R K_c)$, one obtains
\begin{equation}
    V_{k+1} - V_k \leq -\|e_k^z\|_{Q+K_c^\top R K_c}^2 + 2\varepsilon \|P\|\,\|A_{\mathrm{cl}}\|\,\|e_k^z\| + \varepsilon^2 \|P\|.
\end{equation}
This is negative when $\|e_k^z\|$ exceeds a threshold proportional to $\varepsilon$, establishing ultimate boundedness. The explicit bound \eqref{eq:ub_bound} follows from iterating the linear recurrence \eqref{eq:cl_error_dyn} with $\|\delta_k\| \leq 2\varepsilon$ and $\|A_{\mathrm{cl}}\| < 1$: unrolling $j$ steps from any $k \geq k_0$ gives
\begin{equation}
    \|e_{k+j}^z\| \leq \|A_{\mathrm{cl}}\|^j \|e_k^z\| + \frac{1 - \|A_{\mathrm{cl}}\|^j}{1 - \|A_{\mathrm{cl}}\|}\cdot 2\varepsilon,
\end{equation}
and taking $j\to\infty$ yields the stated bound.
\end{proof}

\subsubsection*{Physical-state tracking bound}

Since $e_k^x = Ce_k^z$ with $C=[I_n\;0]$ and $\|C\|=1$, Proposition~\ref{prop:ub} directly implies
\begin{equation}
    \|x_k - x_k^{\mathrm{ref}}\| = \|C e_k^z\| \leq \|e_k^z\| \leq \frac{2\varepsilon}{1 - \|A_{\mathrm{cl}}\|}
    \label{eq:phys_bound}
\end{equation}
for all $k \geq k_0$. Hence the hierarchical controller imposes no penalty on physical-state tracking precision: the bound for $\|x_k - x_k^{\mathrm{ref}}\|$ is identical to that for the full lifted error $\|e_k^z\|$.

Proposition~\ref{prop:ub} and bound~\eqref{eq:phys_bound} reveal a direct connection between model quality and closed-loop performance. The ultimate bound grows with $\varepsilon$ and shrinks as $\|A_{\mathrm{cl}}\|$ decreases, i.e., as the closed-loop eigenvalues are pushed further inside the unit disk. This motivates two practical guidelines.

First, improving the lifting map directly reduces $\varepsilon$ and thereby tightens the tracking bound. The experimental results in Section~\ref{sec:experiments} show that the proposed temporal lifting achieves substantially lower one-step prediction error than baseline methods, which implies a smaller $\varepsilon$ and, via~\eqref{eq:phys_bound}, a tighter theoretical performance guarantee.

Second, the LQR weights $Q_x$, $Q_h$, $\varepsilon$, and $R$ jointly influence $\|A_{\mathrm{cl}}\|$: increasing $Q_x$ relative to $R$ drives the physical-state closed-loop eigenvalues further inside the unit disk, while the small factor $\varepsilon$ attached to $Q_h$ ensures that memory-state feedback remains an auxiliary correction. In practice, the weight selection balances tracking tightness against actuator constraints.

\section{Experiments}
\label{sec:experiments}

The experimental study consists of completed trajectory-prediction comparison and ablation analysis, followed by a brief description of the intended trajectory-tracking validation setting. The reported evidence in this paper is therefore limited to model-identification quality under open-loop evaluation.

\subsection{Experimental Setup}

\subsubsection{Platform-Specific Datasets}

The present study uses two platform-specific controlled-system datasets with complementary roles.

\textbf{Platform A: manipulator dataset.} This dataset corresponds to a 6-state, 3-input controlled mechanical system. The available split contains 15 training trajectories, 4 validation trajectories, and 1 test trajectory, with 74,655, 19,909, and 4,976 recorded time steps, respectively.

\textbf{Platform B: two-state dataset.} This dataset corresponds to a second controlled setting with state dimension $n=2$ and control dimension $m=2$. The available split contains 49 training trajectories, 21 validation trajectories, and 1 test trajectory, with 24,150, 10,351, and 500 recorded time steps, respectively.

In both datasets, valid samples are constructed from history windows that do not cross trajectory boundaries. State and control normalization statistics are computed on the training split and then reused for validation and test data. The two datasets do not yet form a perfectly uniform benchmark; instead, they play complementary roles in the present paper. Platform A provides the most complete cross-model rollout comparison, whereas Platform B provides the broader hyperparameter sweep used for ablation.

\subsubsection{Compared Methods}

The comparative study is restricted to Koopman model families supported by the present study and evaluated under a common testing pipeline:
\begin{itemize}
    \item \textbf{Traditional EDMD}, which uses hand-crafted lifted observables.
    \item \textbf{MLP-Koopman}, which learns an instantaneous-state lifting map with an MLP encoder.
    \item \textbf{Transformer-Koopman}, the proposed temporal lifting model based on history windows and patch-based Transformer encoding.
\end{itemize}

\subsubsection{Evaluation Protocol}

All methods are evaluated under the same open-loop autoregressive rollout protocol described in Section~\ref{sec:rollout}. Each model is initialized with a real history window from the test trajectory and recursively feeds back its own predicted state as in \eqref{eq:rollout_step}. The main metrics are rollout RMSE and MAE over the predicted trajectory.

For the ablation study, all reported configurations on Platform B were trained and tested under the same pipeline. The resulting normalized one-step test RMSE values therefore provide a controlled comparison of the effect of architectural choices within that dataset.

\subsection{Trajectory Prediction Accuracy Comparison}

The most complete cross-model rollout comparison is currently available on Platform A. Table~\ref{tab:prediction_platform_a} reports the trajectory-level results.

\begin{table}[!t]
\renewcommand{\arraystretch}{1.2}
\caption{Open-loop trajectory prediction results on Platform A}
\label{tab:prediction_platform_a}
\centering
\begin{tabular}{lcc}
\hline\hline
Method & RMSE & MAE \\
\hline
Traditional EDMD & 2.6206 & 1.9544 \\
MLP-Koopman & 0.5074 & 0.2511 \\
Transformer-Koopman & \textbf{0.1768} & \textbf{0.0917} \\
\hline\hline
\end{tabular}
\end{table}

The results in Table~\ref{tab:prediction_platform_a} show a clear pattern. Moving from hand-crafted observables to learned liftings already reduces rollout error substantially, as seen in the gap between traditional EDMD and MLP-Koopman. Introducing finite-history temporal lifting yields a further reduction, indicating that the improvement is not only due to using a neural encoder, but also to how the lifting map is constructed. In the present setting, explicitly encoding short trajectory context produces a materially more accurate Koopman model than instantaneous-state lifting alone.

These results should still be interpreted with the correct scope. Platform A is used as the main cross-model case because it is the dataset for which the available evidence provides the most complete like-for-like rollout comparison. Even so, the result is directly relevant to the paper's modeling objective, since recursive rollout accuracy is the most practical test of whether the learned lifted dynamics are useful.

\subsection{Ablation Study}

The ablation study is conducted on Platform B using the available configuration sweep of the proposed model. The best-performing configuration uses history length $P=6$, patch length $p=2$, latent dimension $d=8$, Transformer dimension $d_{\mathrm{model}}=16$, three Transformer layers, two attention heads, and EDMD regularization $\lambda=10^{-4}$.

Table~\ref{tab:ablation_platform_b} reports representative ablation results.

\begin{table}[!t]
\renewcommand{\arraystretch}{1.2}
\caption{Representative ablation results on Platform B}
\label{tab:ablation_platform_b}
\centering
\begin{tabular}{>{\raggedright\arraybackslash}p{0.62\columnwidth}c}
\hline\hline
Variant & Test RMSE \\
\hline
Best setting ($P=6$, $p=2$, 3 layers, 2 heads) & \textbf{0.01408} \\
Reduce history to $P=4$ & 0.01525 \\
Reduce history to $P=2$ & 0.01710 \\
Degenerate patching ($P=2$, $p=1$) & 0.03362 \\
Use 1 layer and 1 head & 0.01513 \\
Use larger latent/model size ($d=16$, $d_{\mathrm{model}}=32$) & 0.01481 \\
Increase regularization to $\lambda=1$ & 0.01591 \\
Increase regularization to $\lambda=10$ & 0.02153 \\
\hline\hline
\end{tabular}
\\[2pt]
{\footnotesize Normalized one-step test RMSE; lower is better.}
\end{table}

Several trends are consistent across these ablations. First, reducing the history length degrades performance, which supports the premise that short-term temporal information is useful for constructing the lifting map. Second, the large error increase in the degenerate short-history setting shows that the gain is not explained by parameter count alone; the patch structure itself matters. Third, increasing latent and embedding size does not automatically improve performance, suggesting that compact temporal liftings can already be sufficient for this dataset. Finally, overly strong EDMD regularization is detrimental, whereas moderate regularization provides a better trade-off.

From a modeling perspective, these results are important because they connect the architecture to the observed gains. The method benefits from three ingredients working together: enough history to capture local dynamics, patch-based compression before attention, and a regularized latent linear fit that avoids over-conditioning the operator regression.

\subsection{Planned Trajectory-Tracking Validation}

The final component of the intended validation pipeline is controller-based trajectory tracking. Since the present study is centered on offline model learning and open-loop evaluation, the manuscript reports the validation protocol rather than completed closed-loop numbers.

The planned comparison is conceptually simple. The controller structure will be kept fixed, and only the learned Koopman model will be changed. In this way, any change in tracking performance can be attributed primarily to model quality rather than to controller redesign. The candidate models for this comparison are Traditional EDMD, MLP-Koopman, and Transformer-Koopman.

The intended evaluation metrics include tracking RMSE, peak tracking error, and control effort or input smoothness. This protocol is consistent with the control-oriented formulation in Section~\ref{sec:control} and provides a direct path toward a later closed-loop study once the controller pipeline is finalized.

\section{Conclusion}
\label{sec:conclusion}

This paper presented Transformer-Koopman, a temporal lifting framework for learning Koopman models of controlled nonlinear systems from finite history windows. The method uses non-overlapping temporal patches and a Transformer encoder to build a compact history-dependent lifting map, propagates the lifted state with linear controlled dynamics, and preserves exact physical-state recovery through a fixed projection. In this way, the resulting model is not only predictive but also structurally suitable for control-oriented evaluation and later controller design.

The completed experiments support the main identification claim of the paper. On the manipulator dataset, the proposed method achieves substantially lower rollout error than the implemented traditional EDMD and MLP-Koopman baselines. On the second dataset, the ablation study shows that temporal context, patch design, and moderate regularization all contribute meaningfully to model quality. Together, these results indicate that a more accurate Koopman predictor can be obtained when the lifting map is learned from temporally structured inputs rather than from the current state alone.

Closed-loop tracking validation is not yet included in the present manuscript. The most direct next step is therefore to compare the learned predictors under a common controller structure and assess whether improved predictive accuracy also translates into improved control performance. A second natural extension is to study richer lifting inputs, including joint state-input histories, within the same control-oriented Koopman framework.

\begin{thebibliography}{99}

\bibitem{koopman1931hamiltonian}
B. O. Koopman, ``Hamiltonian systems and transformation in Hilbert space,'' \emph{Proc. Natl. Acad. Sci. USA}, vol. 17, no. 5, pp. 315--318, 1931.

\bibitem{koopman1932spectra}
B. O. Koopman and J. von Neumann, ``Dynamical systems of continuous spectra,'' \emph{Proc. Natl. Acad. Sci. USA}, vol. 18, no. 3, pp. 255--263, 1932.

\bibitem{mezic2005spectral}
I. Mezi\'c, ``Spectral properties of dynamical systems, model reduction and decompositions,'' \emph{Nonlinear Dyn.}, vol. 41, pp. 309--325, 2005.

\bibitem{ljung1999system}
L. Ljung, \emph{System Identification: Theory for the User}, 2nd ed. Upper Saddle River, NJ, USA: Prentice Hall, 1999.

\bibitem{vanoverschee1994n4sid}
P. Van Overschee and B. De Moor, ``N4SID: Subspace algorithms for the identification of combined deterministic-stochastic systems,'' \emph{Automatica}, vol. 30, no. 1, pp. 75--93, 1994.

\bibitem{nelles2020nonlinear}
O. Nelles, \emph{Nonlinear System Identification: From Classical Approaches to Neural Networks and Fuzzy Models}, 2nd ed. Cham, Switzerland: Springer, 2020.

\bibitem{rowley2009spectral}
C. W. Rowley, I. Mezi\'c, S. Bagheri, P. Schlatter, and D. S. Henningson, ``Spectral analysis of nonlinear flows,'' \emph{J. Fluid Mech.}, vol. 641, pp. 115--127, 2009.

\bibitem{williams2015data}
M. O. Williams, I. G. Kevrekidis, and C. W. Rowley, ``A data--driven approximation of the Koopman operator: Extending dynamic mode decomposition,'' \emph{J. Nonlinear Sci.}, vol. 25, no. 6, pp. 1307--1346, 2015.

\bibitem{williams2015kernel}
M. O. Williams, C. W. Rowley, and I. G. Kevrekidis, ``A kernel-based method for data-driven Koopman spectral analysis,'' \emph{J. Comput. Dyn.}, vol. 2, no. 2, pp. 247--265, 2015.

\bibitem{arbabi2017ergodic}
H. Arbabi and I. Mezi\'c, ``Ergodic theory, dynamic mode decomposition, and computation of spectral properties of the Koopman operator,'' \emph{SIAM J. Appl. Dyn. Syst.}, vol. 16, no. 4, pp. 2096--2126, 2017.

\bibitem{brunton2022modern}
S. L. Brunton, M. Budi\v{s}i\'c, E. Kaiser, and J. N. Kutz, ``Modern Koopman theory for dynamical systems,'' \emph{SIAM Rev.}, vol. 64, no. 2, pp. 229--340, 2022.

\bibitem{otto2021koopman}
S. E. Otto and C. W. Rowley, ``Koopman operators for estimation and control of dynamical systems,'' \emph{Annu. Rev. Control Robot. Auton. Syst.}, vol. 4, pp. 59--87, 2021.

\bibitem{takeishi2017koopman}
N. Takeishi, Y. Kawahara, and T. Yairi, ``Learning Koopman invariant subspaces for dynamic mode decomposition,'' in \emph{Proc. NeurIPS}, 2017, pp. 1130--1140.

\bibitem{yeung2019learning}
E. Yeung, S. Kundu, and N. Hodas, ``Learning deep neural network representations for Koopman operators of nonlinear dynamical systems,'' in \emph{Proc. ACC}, 2019, pp. 4832--4839.

\bibitem{lusch2018deep}
B. Lusch, J. N. Kutz, and S. L. Brunton, ``Deep learning for universal linear embeddings of nonlinear dynamics,'' \emph{Nat. Commun.}, vol. 9, no. 1, p. 4950, 2018.

\bibitem{morton2019deep}
J. Morton, F. D. Witherden, and M. J. Kochenderfer, ``Deep variational Koopman models: Inferring Koopman observations for uncertainty-aware dynamics modeling and control,'' in \emph{Proc. IJCAI}, 2019, pp. 5113--5119.

\bibitem{azencot2020forecasting}
O. Azencot, N. B. Erichson, V. Lin, and M. W. Mahoney, ``Forecasting sequential data using consistent Koopman autoencoders,'' in \emph{Proc. ICML}, 2020, pp. 475--485.

\bibitem{yuan2022variational}
J. Yuan, A. A. Heydari, and Z. Zhang, ``Variational Koopman models: Slow coherent structures and uncertainty quantification,'' \emph{IEEE Trans. Autom. Control}, vol. 67, no. 11, pp. 5905--5920, 2022.

\bibitem{otto2019linearly}
S. E. Otto and C. W. Rowley, ``Linearly-recurrent autoencoder networks for learning dynamics,'' \emph{SIAM J. Appl. Dyn. Syst.}, vol. 18, no. 1, pp. 558--593, 2019.

\bibitem{rank2019recurrent}
M. Rank, A. G. Georgescu, and S. Kirchner, ``Recurrent Koopman operator for data-driven modeling of nonlinear dynamics,'' \emph{IEEE Control Syst. Lett.}, vol. 3, no. 4, pp. 912--917, 2019.

\bibitem{forgione2021dynonet}
M. Forgione and D. Varagnolo, ``DynoNet: A neural network architecture for learning dynamical systems,'' \emph{Int. J. Adapt. Control Signal Process.}, vol. 35, no. 4, pp. 612--629, 2021.

\bibitem{brunton2016sindy}
S. L. Brunton, J. L. Proctor, and J. N. Kutz, ``Discovering governing equations from data by sparse identification of nonlinear dynamical systems,'' \emph{Proc. Natl. Acad. Sci. USA}, vol. 113, no. 15, pp. 3932--3937, 2016.

\bibitem{proctor2016dmdc}
J. L. Proctor, S. L. Brunton, and J. N. Kutz, ``Dynamic mode decomposition with control,'' \emph{SIAM J. Appl. Dyn. Syst.}, vol. 15, no. 1, pp. 142--161, 2016.

\bibitem{proctor2018generalizing}
J. L. Proctor, S. L. Brunton, and J. N. Kutz, ``Generalizing Koopman theory to allow for inputs and control,'' \emph{SIAM J. Appl. Dyn. Syst.}, vol. 17, no. 1, pp. 909--930, 2018.

\bibitem{brunton2016invariant}
S. L. Brunton, B. W. Brunton, J. L. Proctor, and J. N. Kutz, ``Koopman invariant subspaces and finite linear representations of nonlinear dynamical systems for control,'' \emph{PLOS ONE}, vol. 11, no. 2, e0150171, 2016.

\bibitem{korda2018linear}
M. Korda and I. Mezi\'c, ``Linear predictors for nonlinear dynamical systems: Koopman operator meets model predictive control,'' \emph{Automatica}, vol. 93, pp. 149--160, 2018.

\bibitem{han2020deep}
Y. Han, W. Hao, and U. Vaidya, ``Deep learning of Koopman representation for control,'' arXiv:2010.07546, 2020.

\bibitem{shi2022deep}
H. Shi and M. Q.-H. Meng, ``Deep Koopman operator with control for nonlinear systems,'' \emph{IEEE Robot. Autom. Lett.}, vol. 7, no. 3, pp. 7700--7707, 2022.

\bibitem{strasser2026overview}
R. Str\"asser, K. Worthmann, I. Mezi\'c, J. Berberich, M. Schaller, and F. Allg\"ower, ``An overview of Koopman-based control: From error bounds to closed-loop guarantees,'' \emph{Annu. Rev. Control}, vol. 61, art. 101035, 2026.

\bibitem{duran2025control}
J. F. Dur\'an-Sig\"uenza, L. I. Minchala, L. E. Garza-Casta\~n\'on, and H. Zhang, ``Control based on the Koopman operator: A comprehensive review,'' \emph{J. Franklin Inst.}, vol. 362, no. 18, art. 108256, 2025.

\bibitem{klus2018data}
S. Klus, P. Koltai, and C. Sch\"utte, ``Data-driven model reduction and transfer operator approximation,'' \emph{J. Nonlinear Sci.}, vol. 28, no. 3, pp. 985--1010, 2018.

\bibitem{vaswani2017attention}
A. Vaswani \emph{et al.}, ``Attention is all you need,'' in \emph{Proc. NeurIPS}, 2017, pp. 5998--6008.

\bibitem{zhou2021informer}
H. Zhou \emph{et al.}, ``Informer: Beyond efficient Transformer for long sequence time-series forecasting,'' in \emph{Proc. AAAI}, 2021, pp. 11106--11115.

\bibitem{wu2021autoformer}
H. Wu, J. Xu, J. Wang, and M. Long, ``Autoformer: Decomposition Transformers with auto-correlation for long-term series forecasting,'' in \emph{Proc. NeurIPS}, 2021, pp. 22419--22430.

\bibitem{zhou2022fedformer}
T. Zhou \emph{et al.}, ``FEDformer: Frequency enhanced decomposed Transformer for long-term series forecasting,'' in \emph{Proc. ICML}, 2022, pp. 27268--27286.

\bibitem{nie2023time}
Y. Nie, N. H. Nguyen, P. Sinthong, and J. Kalagnanam, ``A time series is worth 64 words: Long-term forecasting with Transformers,'' in \emph{Proc. ICLR}, 2023.

\bibitem{zeng2023transformers}
A. Zeng, M. Chen, L. Zhang, and Q. Xu, ``Are Transformers effective for time series forecasting?,'' in \emph{Proc. AAAI}, 2023, pp. 11121--11128.

\bibitem{liu2024itransformer}
Y. Liu \emph{et al.}, ``iTransformer: Inverted Transformers are effective for time series forecasting,'' in \emph{Proc. ICLR}, 2024.

\bibitem{forootani2026compact}
A. Forootani and M. Khosravi, ``Compact transformer variants for synthetic time series forecasting,'' \emph{Neurocomputing}, vol. 677, art. 133140, 2026.

\bibitem{forootani2025synthetic}
A. Forootani and M. Khosravi, ``Synthetic time series forecasting with Transformer architectures: Extensive simulation benchmarks,'' arXiv:2505.20048, 2025.

\bibitem{forootani2026deepkoopformer}
A. Forootani, M. Khosravi, and M. Barati, ``DeepKoopFormer: A Koopman enhanced Transformer based architecture for time series forecasting,'' \emph{npj Artif. Intell.}, vol. 2, art. 35, 2026.

\bibitem{forootani2026spectral}
A. Forootani and R. Iervolino, ``Learnable Koopman-enhanced Transformer-based time series forecasting with spectral control,'' arXiv:2602.02592, 2026.

\bibitem{wang2020linformer}
S. Wang, B. Z. Li, M. Khabsa, H. Fang, and H. Ma, ``Linformer: Self-attention with linear complexity,'' arXiv:2006.04768, 2020.

\bibitem{katharopoulos2020transformers}
A. Katharopoulos, A. Vyas, N. Pappas, and F. Fleuret, ``Transformers are RNNs: Fast autoregressive Transformers with linear attention,'' in \emph{Proc. ICML}, 2020, pp. 5156--5165.

\bibitem{choromanski2021performer}
K. M. Choromanski \emph{et al.}, ``Rethinking attention with Performers,'' in \emph{Proc. ICLR}, 2021.

\end{thebibliography}

\end{document}